\begin{abstract}
LLM-as-a-Judge research usually asks whether a model assigns the right score or
preference. We study an earlier failure mode: before any score is produced, the
judge may omit evaluation dimensions that human evaluators consider essential.
We call this failure evaluation blind spots. On a hard-gold RubricBench
diagnostic split, a single-model evaluation-criteria policy covers only 36.9\% of
human-gold evaluation dimensions, leaving a mean blind-spot rate of 63.1\%
(95\% CI: 57.9--68.2). To make this phenomenon measurable and optimizable, we
introduce Blind-Spot Coverage (BSC), a semantic coverage metric over
human-gold evaluation dimensions with explicit penalties for redundancy and
invalid or hallucinated criteria. BSC turns open-ended criteria elicitation into a
verifiable reward: coverage credit requires semantic match to hidden evaluation
dimensions, while redundancy and fail-closed validity checks penalize duplicate,
invalid, undecidable, irrelevant, or hallucinated criteria. We
instantiate this reward in a two-stage training framework, BlindSpot-RL, that
first constructs proxy-gold criteria through multi-teacher criteria elicitation
and fail-closed verification, then tests GRPO/RLVR optimization of an
evaluation-criteria policy under the BSC reward. The contribution is to cast
open-ended criteria elicitation as verifiable semantic reward
optimization, rather than as surface-form optimization.
The evidence-gated experiments then test when the artifacts are sufficient to
write a dimension-level coverage-change statement about the criteria LLM judges
consider; any paper-facing dimension-recovery wording is permitted only after
the hard-gold, downstream, ablation, dimension-transition, and human-audit gates
support it.
\end{abstract}
